\section{Case Study: Cabin}


\begin{frame}{Project Cabin}
    \mycomment{
    Another example: Harald has built this cabin in his garden from materials left over from his house construction site. He did this without much of a plan, though.

    You can see that this hut was built without an architect: there is no space between the window and the door. There is also no stairs leading up to the hut. There is no water or electricity supply; the use of an extension chord was not a solution.

    But worst of all, after construction he was forced to tear it all down again, because for such a permanent structure you need building permission!
    }

    \showmybitmap{../shared/img/fitting-stairs2.JPG}
\end{frame}

% \begin{frame}{Project Cabin}
%     \showmybitmap{../shared/img/fitting-stairs.JPG}
% \end{frame}

\begin{frame}{Project Cabin}
    \showmytikz{../shared/tikz/tikz-PM-gantt-cabin-v1}
    \mycomment{
    This is a sketch for a Gantt diagram for the \enquote{projec cabin}. Doing one thing after the other, the project could easily take up to nine months.    
    }
\end{frame}

\begin{frame}{Project Cabin}
    \showmytikz{../shared/tikz/tikz-PM-gantt-cabin-v2}
    \mycomment{
    This amends the plan by procurement. At each stage, Harald had to notice that he could not complete the project with the material in place. He had to procure (more) wood, he had to buy sockets and other utility fittings, and he also had to equip the interior of the hut with curtains, furniture and so forth.


    So just by forgetting that procurement of requirements takes time, the project time has stretched out dramatically.
    
    }
\end{frame}

\begin{frame}{Project Cabin}
    \showmytikz{../shared/tikz/tikz-PM-gantt-cabin-v3}
    \mycomment{
    
    But that wasn't all. He actually also had to procure (find) the workers to construct the hut---and that was something that seriously delayed everything.

    He needed different contractors for the different types of work: the carpenters to build the cabin; the utility guys to fit in electrity (and water); and someone to take care of the interior, such as painting, decoration and furnishing.

    The problem is, of course, that here the utility guys can only start working once the cabin is in place, and so forth.
    One trade waits for the other.
    
    }
\end{frame}

\begin{frame}{Project Cabin}
    \showmytikz{../shared/tikz/tikz-PM-gantt-cabin-v4}
    \mycomment{
    If you look closer, you could have the ultility guys and the interior guys work in parallel to some extent. And they could already start procuring their requirement before the carpenters have topped up the building.

    Now it becomes visible that in month 5, all three trades work at the same time. Can they really? And if they do, where do they sleep? Who is coordinating the site?
    
    }
\end{frame}

\begin{frame}{Project Cabin}
    \mycomment{
    So you will always have someone working on the project all the time: the central management, and the site manager. This can be the same person---in our case, Harald. So while different people are working on the cabin, he will have to do work all the time---just the management an planning will keep him quite busy.
    }
    \showmytikz{../shared/tikz/tikz-PM-gantt-cabin-v5}
    \end{frame}

\begin{frame}{Project Cabin}
    
    \mytext{
    
    \begin{itemize}
        \item Planning is super important
        \item Planning itself can and must be planned
        \item Plans may change!
        \begin{itemize}
            \item Make flexible plans!
            \item Plan for alternatives!
        \end{itemize}
        \item Planning looks at \textbf{interfaces}
        \item Planning is communication
    \end{itemize}
    
    }
\end{frame}

 
